% Options for packages loaded elsewhere
\PassOptionsToPackage{unicode}{hyperref}
\PassOptionsToPackage{hyphens}{url}
%
\documentclass[
  12pt,
]{article}
\usepackage{amsmath,amssymb}
\usepackage{iftex}
\ifPDFTeX
  \usepackage[T1]{fontenc}
  \usepackage[utf8]{inputenc}
  \usepackage{textcomp} % provide euro and other symbols
\else % if luatex or xetex
  \usepackage{unicode-math} % this also loads fontspec
  \defaultfontfeatures{Scale=MatchLowercase}
  \defaultfontfeatures[\rmfamily]{Ligatures=TeX,Scale=1}
\fi
\usepackage{lmodern}
\ifPDFTeX\else
  % xetex/luatex font selection
\fi
% Use upquote if available, for straight quotes in verbatim environments
\IfFileExists{upquote.sty}{\usepackage{upquote}}{}
\IfFileExists{microtype.sty}{% use microtype if available
  \usepackage[]{microtype}
  \UseMicrotypeSet[protrusion]{basicmath} % disable protrusion for tt fonts
}{}
\makeatletter
\@ifundefined{KOMAClassName}{% if non-KOMA class
  \IfFileExists{parskip.sty}{%
    \usepackage{parskip}
  }{% else
    \setlength{\parindent}{0pt}
    \setlength{\parskip}{6pt plus 2pt minus 1pt}}
}{% if KOMA class
  \KOMAoptions{parskip=half}}
\makeatother
\usepackage{xcolor}
\usepackage[margin=2.5cm]{geometry}
\usepackage{longtable,booktabs,array}
\usepackage{calc} % for calculating minipage widths
% Correct order of tables after \paragraph or \subparagraph
\usepackage{etoolbox}
\makeatletter
\patchcmd\longtable{\par}{\if@noskipsec\mbox{}\fi\par}{}{}
\makeatother
% Allow footnotes in longtable head/foot
\IfFileExists{footnotehyper.sty}{\usepackage{footnotehyper}}{\usepackage{footnote}}
\makesavenoteenv{longtable}
\usepackage{graphicx}
\makeatletter
\newsavebox\pandoc@box
\newcommand*\pandocbounded[1]{% scales image to fit in text height/width
  \sbox\pandoc@box{#1}%
  \Gscale@div\@tempa{\textheight}{\dimexpr\ht\pandoc@box+\dp\pandoc@box\relax}%
  \Gscale@div\@tempb{\linewidth}{\wd\pandoc@box}%
  \ifdim\@tempb\p@<\@tempa\p@\let\@tempa\@tempb\fi% select the smaller of both
  \ifdim\@tempa\p@<\p@\scalebox{\@tempa}{\usebox\pandoc@box}%
  \else\usebox{\pandoc@box}%
  \fi%
}
% Set default figure placement to htbp
\def\fps@figure{htbp}
\makeatother
\setlength{\emergencystretch}{3em} % prevent overfull lines
\providecommand{\tightlist}{%
  \setlength{\itemsep}{0pt}\setlength{\parskip}{0pt}}
\setcounter{secnumdepth}{-\maxdimen} % remove section numbering
\usepackage{booktabs}
\usepackage{longtable}
\usepackage{array}
\usepackage{multirow}
\usepackage{wrapfig}
\usepackage{float}
\usepackage{colortbl}
\usepackage{pdflscape}
\usepackage{tabu}
\usepackage{threeparttable}
\usepackage{threeparttablex}
\usepackage[normalem]{ulem}
\usepackage{makecell}
\usepackage{xcolor}
\usepackage{bookmark}
\IfFileExists{xurl.sty}{\usepackage{xurl}}{} % add URL line breaks if available
\urlstyle{same}
\hypersetup{
  pdftitle={Trabalho 1 - ME613 : Analise de Regressao},
  hidelinks,
  pdfcreator={LaTeX via pandoc}}

\title{Trabalho 1 - ME613 : Analise de Regressao}
\author{\textbf{Daniela Cielo}: RA: 298825\\
\textbf{Geovani Ariel}: RA: 236124\\
\textbf{Lucas Lima}: RA: 263267}
\date{2º semestre de 2025}

\begin{document}
\maketitle

{
\setcounter{tocdepth}{2}
\tableofcontents
}
\section{Introducao e Motivacao}\label{introducao-e-motivacao}

A expectativa de vida e um indicador fundamental para avaliar o nivel de
saude e desenvolvimento de uma populacao, refletindo nao apenas as
condicoes medicas, mas tambem fatores socioeconomicos, ambientais e de
politicas publicas. Ao sintetizar o tempo medio de vida que uma coorte
hipotetica de individuos teria, caso permanecessem expostos as taxas de
mortalidade observadas em determinado periodo, esse indicador fornece
uma visao abrangente das condicoes de vida em diferentes paises e
contextos historicos.

Nas ultimas decadas, o aumento da expectativa de vida tem sido um dos
fenomenos demograficos mais marcantes em grande parte do mundo,
especialmente em paises desenvolvidos. No entanto, essa evolucao nao e
homogenea: muitas regioes ainda enfrentam desafios importantes
relacionados a mortalidade infantil, doencas infecciosas, desnutricao,
desigualdade de renda e acesso limitado a servicos basicos de saude e
educacao. Essas disparidades revelam que ganhos agregados em expectativa
de vida podem ocultar heterogeneidades relevantes entre paises, grupos
populacionais e periodos de tempo.

Do ponto de vista de politicas publicas, compreender os determinantes da
expectativa de vida e essencial para orientar decisoes estrategicas,
como priorizacao de investimentos em saude, educacao e infraestrutura
social. Alem disso, a analise conjunta de variaveis de mortalidade,
indicadores economicos e fatores sociais permite investigar em que
medida os avancos em desenvolvimento humano se traduzem, de fato, em
melhores condicoes de sobrevivencia e bem-estar da populacao.

Nesse contexto, este trabalho analisa a relacao entre a expectativa de
vida e um conjunto de variaveis explicativas associadas a fatores de
mortalidade, imunizacao, condicoes economicas e fatores sociais,
utilizando um conjunto de dados em painel com paises observados ao longo
do tempo. A partir de tecnicas de analise exploratoria e modelagem de
regressao, busca-se identificar quais fatores apresentam maior
associacao com a expectativa de vida e discutir implicacoes para
politicas publicas e planejamento em saude.

\section{Materiais e Metodos}\label{materiais-e-metodos}

Nesta seção descrevemos o banco de dados utilizado, as variaveis
consideradas e os principais procedimentos adotados no tratamento e na
preparacao dos dados para analise.

\subsection{Base de dados e variaveis}\label{base-de-dados-e-variaveis}

O conjunto de dados utilizado contem informacoes anuais sobre diversos
paises, incluindo:

\begin{itemize}
\tightlist
\item
  \textbf{Expectativa\_Vida}: variavel resposta, medida em anos;
\item
  indicadores de \textbf{mortalidade} (Mortalidade\_Adulto,
  Mortalidade\_Infantil, Mortes\_Menores\_Cinco\_Anos);
\item
  indicadores de \textbf{saude e imunizacao} (Hepatite\_B, Poliomielite,
  Difteria, HIV\_AIDS);
\item
  variaveis \textbf{economicas e sociais} (PIB,
  Composicao\_Renda\_Recursos, Escolaridade, Percentual\_Despesa,
  Despesa\_Total);
\item
  variaveis \textbf{nutricionais e demograficas} (IMC,
  Magreza\_1\_19\_Anos, Magreza\_5\_9\_Anos, Populacao).
\end{itemize}

\begin{longtable}[t]{rrrr}
\caption{\label{tab:unnamed-chunk-3}Resumo geral da base de dados}\\
\toprule
n\_obs & n\_paises & ano\_min & ano\_max\\
\midrule
2938 & 193 & 2000 & 2015\\
\bottomrule
\end{longtable}

A Tabela 1 ilustra as primeiras observacoes do conjunto de dados ja
tratado, incluindo a variavel resposta \textbf{Expectativa\_Vida} e
algumas variaveis explicativas de interesse (Mortalidade\_Adulto,
Escolaridade e Populacao).

\subsection{Avaliacao de valores
faltantes}\label{avaliacao-de-valores-faltantes}

Antes da modelagem, e necessario avaliar a presenca de \textbf{valores
faltantes (NA)}, pois variaveis com muitos NA podem reduzir bastante o
numero de observacoes utilizaveis e comprometer a estabilidade das
estimativas.

\begin{longtable}[t]{lrr}
\caption{\label{tab:unnamed-chunk-4}Quantidade e percentual de valores faltantes por variavel}\\
\toprule
Variavel & Quantidade\_NA & Percentual\_NA\\
\midrule
Populacao & 652 & 22.19\\
Hepatite\_B & 553 & 18.82\\
PIB & 448 & 15.25\\
Despesa\_Total & 226 & 7.69\\
Alcool & 194 & 6.60\\
\addlinespace
Composicao\_Renda\_Recursos & 167 & 5.68\\
Escolaridade & 163 & 5.55\\
IMC & 34 & 1.16\\
Magreza\_1\_19\_Anos & 34 & 1.16\\
Magreza\_5\_9\_Anos & 34 & 1.16\\
\addlinespace
Poliomielite & 19 & 0.65\\
Difteria & 19 & 0.65\\
Expectativa\_Vida & 10 & 0.34\\
Mortalidade\_Adulto & 10 & 0.34\\
Pais & 0 & 0.00\\
\addlinespace
Ano & 0 & 0.00\\
Status\_Pais & 0 & 0.00\\
Mortalidade\_Infantil & 0 & 0.00\\
Percentual\_Despesa & 0 & 0.00\\
Sarampo & 0 & 0.00\\
\addlinespace
Mortes\_Menores\_Cinco\_Anos & 0 & 0.00\\
HIV\_AIDS & 0 & 0.00\\
\bottomrule
\end{longtable}

A Tabela de valores faltantes é usada para nos ajudar decidir quais
variaveis devem ser mantidas ou excluidas da base destinada a modelagem,
em combinacao com informacoes de correlacao apresentadas adiante.

\subsection{Estrategia de análise}\label{estrategia-de-anuxe1lise}

A estrategia de análise utilizada neste trabalho segue os seguintes
passos:

\begin{enumerate}
\def\labelenumi{\arabic{enumi}.}
\tightlist
\item
  \textbf{Análise descritiva da variavel resposta}, incluindo resumo
  estatistico e comparacao por Status\_Pais e ao longo do tempo;
\item
  \textbf{Análise de correlacao}, avaliando:

  \begin{itemize}
  \tightlist
  \item
    correlacao entre cada variavel numerica e a Expectativa\_Vida;
  \item
    correlacao entre os proprios preditores, para identificar possivel
    multicolinearidade;
  \end{itemize}
\item
  \textbf{Definicao da base para modelagem}, com:

  \begin{itemize}
  \tightlist
  \item
    exclusao de variaveis com alta proporcao de NA e baixa correlacao
    com a resposta (Populacao e Hepatite\_B);
  \item
    combinacao de variaveis altamente correlacionadas e conceitualmente
    proximas (Magreza\_1\_19\_Anos e Magreza\_5\_9\_Anos) em uma unica
    variavel \textbf{Magreza\_media}.
  \end{itemize}
\end{enumerate}

A modelagem de regressao (na etapa seguinte do trabalho) sera realizada
utilizando essa base reduzida, mais adequada em termos de qualidade da
informacao e estabilidade numerica.

\subsection{Analise da variavel
resposta}\label{analise-da-variavel-resposta}

\begin{longtable}[t]{rrrrrrrr}
\caption{\label{tab:unnamed-chunk-5}Resumo descritivo da variavel resposta}\\
\toprule
n & media & sd & min & q1 & mediana & q3 & max\\
\midrule
2928 & 69.22 & 9.52 & 36.3 & 63.1 & 72.1 & 75.7 & 89\\
\bottomrule
\end{longtable}

O resumo descritivo indica que a expectativa de vida apresenta ampla
variacao entre observacoes, mas com concentracao em torno de valores
intermediarios a altos, sugerindo heterogeneidade entre contextos mais e
menos desenvolvidos.

\begin{longtable}[t]{lrrrrrrrr}
\caption{\label{tab:unnamed-chunk-6}Resumo da expectativa de vida por status do pais}\\
\toprule
Status\_Pais & n & media & sd & min & q1 & mediana & q3 & max\\
\midrule
Developed & 512 & 79.20 & 3.93 & 69.9 & 76.8 & 79.25 & 81.7 & 89\\
Developing & 2416 & 67.11 & 9.01 & 36.3 & 61.1 & 69.00 & 74.0 & 89\\
\bottomrule
\end{longtable}

Os resultados por \textbf{Status\_Pais} evidenciam que paises
desenvolvidos tendem a apresentar maior expectativa de vida media e
menor variabilidade, ao passo que paises em desenvolvimento exibem media
menor e maior dispersao, coerente com desigualdades em saude.

\begin{center}\includegraphics{Trabalho_ME613_files/figure-latex/unnamed-chunk-7-1} \end{center}

\begin{center}\includegraphics{Trabalho_ME613_files/figure-latex/unnamed-chunk-8-1} \end{center}

Os graficos temporais mostram uma tendencia de aumento da expectativa de
vida ao longo dos anos para ambos os grupos de paises, com nivel
consistentemente mais alto entre os paises desenvolvidos.

\subsection{Estrutura de correlacoes}\label{estrutura-de-correlacoes}

\subsubsection{Correlacao com a variavel
resposta}\label{correlacao-com-a-variavel-resposta}

\begin{longtable}[t]{lr}
\caption{\label{tab:unnamed-chunk-9}Correlacao de cada variavel numerica com a resposta}\\
\toprule
Variavel & Expectativa\_Vida\\
\midrule
Expectativa\_Vida & 1.000\\
Escolaridade & 0.752\\
Composicao\_Renda\_Recursos & 0.725\\
Mortalidade\_Adulto & -0.696\\
IMC & 0.568\\
\addlinespace
HIV\_AIDS & -0.557\\
Difteria & 0.479\\
Magreza\_1\_19\_Anos & -0.477\\
Magreza\_5\_9\_Anos & -0.472\\
Poliomielite & 0.466\\
\addlinespace
PIB & 0.461\\
Alcool & 0.405\\
Percentual\_Despesa & 0.382\\
Hepatite\_B & 0.257\\
Mortes\_Menores\_Cinco\_Anos & -0.223\\
\addlinespace
Despesa\_Total & 0.218\\
Mortalidade\_Infantil & -0.197\\
Ano & 0.170\\
Sarampo & -0.158\\
Populacao & -0.022\\
\bottomrule
\end{longtable}

A tabela acima organiza as variaveis numericas em ordem decrescente de
correlacao (em valor absoluto) com a Expectativa\_Vida. Como esperado,
indicadores de mortalidade (como Mortalidade\_Adulto), desnutricao
(Magreza\_1\_19\_Anos, Magreza\_5\_9\_Anos) e doencas como HIV\_AIDS
apresentam correlacoes negativas relevantes, enquanto variaveis
socioeconomicas (Escolaridade, Composicao\_Renda\_Recursos, PIB) tendem
a apresentar correlacao positiva com a expectativa de vida.

\subsubsection{Correlacao entre
preditores}\label{correlacao-entre-preditores}

\begin{longtable}[t]{llr}
\caption{\label{tab:unnamed-chunk-10}Pares de variáveis numéricas com correlação forte (|r| >= 0.7)}\\
\toprule
var1 & var2 & cor\\
\midrule
Mortalidade\_Infantil & Mortes\_Menores\_Cinco\_Anos & 0.997\\
Magreza\_1\_19\_Anos & Magreza\_5\_9\_Anos & 0.939\\
Percentual\_Despesa & PIB & 0.899\\
Composicao\_Renda\_Recursos & Escolaridade & 0.800\\
\bottomrule
\end{longtable}

A identificacao de pares de preditores com \textbf{alta correlacao} e
importante para evitar problemas de multicolinearidade na regressao.
Nesse contexto, destaca-se o par formado por
\textbf{Magreza\_1\_19\_Anos} e \textbf{Magreza\_5\_9\_Anos}, que
apresentam correlacao elevada entre si e medem fenomenos muito proximos
(magreza em faixas etarias infantis semelhantes).

\subsubsection{Matriz de correlacao para um subconjunto de
variaveis}\label{matriz-de-correlacao-para-um-subconjunto-de-variaveis}

\begin{center}\includegraphics{Trabalho_ME613_files/figure-latex/unnamed-chunk-11-1} \end{center}

A matriz de correlacao para esse subconjunto reforca a interpretacao de
que:

\begin{itemize}
\tightlist
\item
  a expectativa de vida esta fortemente associada a fatores de
  mortalidade, nutricionais e socioeconomicos;
\item
  Magreza\_1\_19\_Anos e Magreza\_5\_9\_Anos sao fortemente
  correlacionadas entre si, justificando a criacao de um indicador
  sintetico para a fase de modelagem.
\end{itemize}

\subsection{Construcao pratica da base de
modelagem}\label{construcao-pratica-da-base-de-modelagem}

Com base nas informacoes de NA e de correlacao, adotamos duas decisoes
para a construcao da base de modelagem:

\begin{enumerate}
\def\labelenumi{\arabic{enumi}.}
\item
  \textbf{Exclusao de Populacao e Hepatite\_B}: ambas apresentaram
  \textbf{alta proporcao de valores faltantes} e \textbf{correlacao
  fraca com Expectativa\_Vida}, de modo que sua inclusao no modelo
  reduziria o numero de casos disponiveis sem trazer ganho relevante de
  explicacao.
\item
  \textbf{Combinacao de Magreza\_1\_19\_Anos e Magreza\_5\_9\_Anos em
  Magreza\_media}: as duas variaveis sao altamente correlacionadas e
  representam a mesma dimensao (magreza infantil em faixas etarias
  proximas). Para reduzir dimensionalidade e multicolinearidade,
  definimos uma unica variavel:
\end{enumerate}

\(begin:math:display\)
\textbackslash text\{Magreza\textbackslash\_media\} =
\textbackslash frac\{\textbackslash text\{Magreza\textbackslash\_1\textbackslash\_19\textbackslash\_Anos\}
+
\textbackslash text\{Magreza\textbackslash\_5\textbackslash\_9\textbackslash\_Anos\}\}\{2\}.
\(end:math:display\)

Na base de modelagem, utilizamos apenas \textbf{Magreza\_media},
mantendo as variaveis originais apenas para a analise exploratoria.

\begin{longtable}[t]{rr}
\caption{\label{tab:unnamed-chunk-12}Dimensoes da base de dados para modelagem}\\
\toprule
Linhas & Colunas\\
\midrule
2938 & 13\\
\bottomrule
\end{longtable}

\begin{longtable}[t]{l}
\caption{\label{tab:unnamed-chunk-13}Variáveis numéricas na base de modelagem}\\
\toprule
Variável\\
\midrule
Expectativa\_Vida\\
Mortalidade\_Adulto\\
Alcool\\
Percentual\_Despesa\\
IMC\\
\addlinespace
Poliomielite\\
Difteria\\
HIV\_AIDS\\
Composicao\_Renda\_Recursos\\
Escolaridade\\
\addlinespace
Magreza\_media\\
\bottomrule
\end{longtable}

A base \texttt{dados\_modelo} sera utilizada na etapa de regressao,
garantindo melhor equilibrio entre qualidade dos dados, relevancia
estatistica dos preditores e estabilidade numerica.

A análise exploratória permitiu caracterizar a expectativa de vida e
seus possíveis determinantes em diferentes países e anos. De forma
geral:

\begin{itemize}
\tightlist
\item
  a expectativa de vida apresenta variabilidade considerável entre
  países, mas tende a ser maior e menos dispersa em países classificados
  como desenvolvidos;
\item
  observa-se uma tendência clara de aumento da expectativa de vida ao
  longo do período analisado, coerente com avanços globais em saúde e
  condições de vida;
\item
  variáveis relacionadas à mortalidade (especialmente
  \textbf{Mortalidade\_Adulto}), desnutrição
  (\textbf{Magreza\_1\_19\_Anos}, \textbf{Magreza\_5\_9\_Anos}) e
  doenças como \textbf{HIV\_AIDS} estão fortemente associadas a menores
  níveis de expectativa de vida;
\item
  variáveis socioeconômicas, como \textbf{Escolaridade},
  \textbf{Composicao\_Renda\_Recursos} e \textbf{PIB}, exibem correlação
  positiva com a expectativa de vida, sugerindo que desenvolvimento
  humano e econômico está ligado a melhores condições de saúde
  populacional.
\end{itemize}

Do ponto de vista da preparação dos dados para modelagem, a combinação
entre informação de valores faltantes e correlações levou a decisões
importantes:

\begin{itemize}
\tightlist
\item
  exclusão de \textbf{Populacao}, \textbf{Hepatite\_B} e \textbf{PIB} da
  base de modelagem, devido à elevada proporção de valores faltantes
  e/ou fraca associação com a resposta;
\item
  criação da variável \textbf{Magreza\_media}, sintetizando a informação
  de magreza infantil em uma única medida, reduzindo redundância entre
  preditores e contribuindo para maior estabilidade dos futuros modelos
  de regressão.
\end{itemize}

Essas etapas de tratamento e análise exploratória fornecem uma base
sólida para a fase seguinte do trabalho, na qual modelos de regressão
serão ajustados utilizando a base \texttt{dados\_modelo}, permitindo
quantificar a importância relativa de cada fator e avaliar o poder
explicativo conjunto sobre a expectativa de vida.

\begin{longtable}[t]{rrr}
\caption{\label{tab:unnamed-chunk-14}Comparação entre base de modelagem original e base completa sem NA}\\
\toprule
Linhas\_originais & Linhas\_sem\_NA & Colunas\\
\midrule
2938 & 2562 & 11\\
\bottomrule
\end{longtable}

\begin{verbatim}
## 
## Call:
## lm(formula = Expectativa_Vida ~ Mortalidade_Adulto + Alcool + 
##     Percentual_Despesa + IMC + Poliomielite + Difteria + HIV_AIDS + 
##     Composicao_Renda_Recursos + Escolaridade + Magreza_media, 
##     data = dados_modelo_comp)
## 
## Residuals:
##      Min       1Q   Median       3Q      Max 
## -20.8557  -2.1882  -0.0663   2.2949  20.6888 
## 
## Coefficients:
##                             Estimate Std. Error t value Pr(>|t|)    
## (Intercept)                5.095e+01  5.338e-01  95.447  < 2e-16 ***
## Mortalidade_Adulto        -1.661e-02  8.302e-04 -20.005  < 2e-16 ***
## Alcool                    -3.087e-02  2.455e-02  -1.258  0.20868    
## Percentual_Despesa         3.644e-04  4.225e-05   8.625  < 2e-16 ***
## IMC                        3.538e-02  5.301e-03   6.674 3.04e-11 ***
## Poliomielite               2.576e-02  4.746e-03   5.427 6.28e-08 ***
## Difteria                   3.148e-02  4.726e-03   6.662 3.30e-11 ***
## HIV_AIDS                  -4.946e-01  1.745e-02 -28.353  < 2e-16 ***
## Composicao_Renda_Recursos  7.455e+00  6.245e-01  11.938  < 2e-16 ***
## Escolaridade               9.416e-01  4.541e-02  20.734  < 2e-16 ***
## Magreza_media             -6.246e-02  2.231e-02  -2.799  0.00516 ** 
## ---
## Signif. codes:  0 '***' 0.001 '**' 0.01 '*' 0.05 '.' 0.1 ' ' 1
## 
## Residual standard error: 3.97 on 2551 degrees of freedom
## Multiple R-squared:  0.8223, Adjusted R-squared:  0.8216 
## F-statistic:  1180 on 10 and 2551 DF,  p-value: < 2.2e-16
\end{verbatim}

\begin{longtable}[]{@{}lr@{}}
\caption{Fatores de inflacao da variancia (VIF) - modelo
completo}\tabularnewline
\toprule\noalign{}
Variavel & VIF \\
\midrule\noalign{}
\endfirsthead
\toprule\noalign{}
Variavel & VIF \\
\midrule\noalign{}
\endhead
\bottomrule\noalign{}
\endlastfoot
Escolaridade & 3.46 \\
Composicao\_Renda\_Recursos & 2.76 \\
Difteria & 1.95 \\
Poliomielite & 1.93 \\
IMC & 1.78 \\
Mortalidade\_Adulto & 1.73 \\
Magreza\_media & 1.62 \\
Alcool & 1.58 \\
HIV\_AIDS & 1.43 \\
Percentual\_Despesa & 1.29 \\
\end{longtable}

\begin{longtable}[]{@{}rrrrrrr@{}}
\caption{Resumo global do modelo selecionado via AIC
(stepwise)}\tabularnewline
\toprule\noalign{}
r.squared & adj.r.squared & sigma & AIC & BIC & df & df.residual \\
\midrule\noalign{}
\endfirsthead
\toprule\noalign{}
r.squared & adj.r.squared & sigma & AIC & BIC & df & df.residual \\
\midrule\noalign{}
\endhead
\bottomrule\noalign{}
\endlastfoot
0.822 & 0.822 & 3.97 & 14347.82 & 14412.15 & 9 & 2552 \\
\end{longtable}

\begin{longtable}[]{@{}lrrrr@{}}
\caption{Estimativas dos coeficientes do modelo selecionado
(stepwise)}\tabularnewline
\toprule\noalign{}
term & estimate & std.error & statistic & p.value \\
\midrule\noalign{}
\endfirsthead
\toprule\noalign{}
term & estimate & std.error & statistic & p.value \\
\midrule\noalign{}
\endhead
\bottomrule\noalign{}
\endlastfoot
(Intercept) & 50.994 & 0.532 & 95.786 & 0.00 \\
Mortalidade\_Adulto & -0.017 & 0.001 & -20.173 & 0.00 \\
Percentual\_Despesa & 0.000 & 0.000 & 8.534 & 0.00 \\
IMC & 0.036 & 0.005 & 6.726 & 0.00 \\
Poliomielite & 0.026 & 0.005 & 5.435 & 0.00 \\
Difteria & 0.032 & 0.005 & 6.686 & 0.00 \\
HIV\_AIDS & -0.496 & 0.017 & -28.518 & 0.00 \\
Composicao\_Renda\_Recursos & 7.431 & 0.624 & 11.902 & 0.00 \\
Escolaridade & 0.925 & 0.043 & 21.298 & 0.00 \\
Magreza\_media & -0.056 & 0.022 & -2.590 & 0.01 \\
\end{longtable}

A partir do modelo completo, foi realizada uma seleção automática de
variáveis via critério de informação de Akaike (AIC), utilizando um
procedimento \emph{stepwise} bidirecional. O modelo resultante,
denominado \texttt{modelo\_step}, apresenta menor AIC/BIC e R² ajustado
ligeiramente superior, indicando melhor compromisso entre ajuste e
parcimônia.

\begin{center}\includegraphics{Trabalho_ME613_files/figure-latex/unnamed-chunk-17-1} \end{center}

\begin{longtable}[t]{rrrl}
\caption{\label{tab:unnamed-chunk-17}Teste de Breusch–Pagan para homocedasticidade (modelo selecionado)}\\
\toprule
statistic & p.value & parameter & method\\
\midrule
329.973 & 0 & 9 & studentized Breusch-Pagan test\\
\bottomrule
\end{longtable}

\section{Referencias bibliograficas}\label{referencias-bibliograficas}

\begin{itemize}
\tightlist
\item
  World Health Organization (WHO). World Health Statistics.\\
\item
  Montgomery, D. C., Peck, E. A., \& Vining, G. G. (2012).
  \emph{Introduction to Linear Regression Analysis}.\\
\item
  Outros artigos e relatorios relevantes utilizados na motivacao e
  discussao.
\end{itemize}

\end{document}
